\UnnumberedChapter{会议论文集引言}
\SourcePage{8}
这三卷（第100、101、102卷）汇集了1981年7月6日至11日在Luminy举行的“奇异空间上的分析与拓扑”讨论会的论文。会议涉及若干主题，下面作简要介绍。

已知正则完整线性微分\term{complex}{复形}的解\term{sheaf}{层}\term{complex}{复形}具有解析\term{constructible}{可构造}的\term{cohomology}{上同调}。所谓Riemann--Hilbert对应，在\term{constructible}{可构造}\term{complex}{复形}与正则完整\term{complex}{复形}之间建立了\term{equivalence}{等价}。在这一对应下，一个正则\term{holonomic-module}{完整模}对应于一个\term{constructible}{可构造}\term{complex}{复形}，而这个\term{complex}{复形}通常有多个非零\term{cohomology}{上同调}\term{sheaf}{层}。这些特殊的\term{constructible}{可构造}\term{complex}{复形}称为\term{perverse-sheaf}{反常层}，可以用纯粹的拓扑条件刻画。同样的刻画也适用于任意域（例如有限域）上不必\term{smooth}{光滑}的簇上的\term{constructible}{可构造}\term{etale}{平展}\term{sheaf}{层}\term{complex}{复形}。因此便得到\term{etale}{平展}\term{perverse-sheaf}{反常层}的概念。虽然这一\term{complex}{复形}范畴不再对应于微分模，我们仍可证明，它实际上构成了一个新的\term{abelian-category}{阿贝尔范畴}。

\term{etale}{平展}\term{perverse-sheaf}{反常层}的一个重要之处在于：利用通过特殊化从特征零过渡到特征$p$的一般定理，只要能在有限域上证明有关结论，就可以得到复数域上\term{perverse-sheaf}{反常层}的相应结果。\term{hypercohomology}{超上同调}给出\term{intersection-homology}{交叉同调}的\term{intersection-complex}{交叉复形}是一个\term{perverse-sheaf}{反常层}，它在有限域上特殊化为\term{etale}{平展}\term{intersection-complex}{交叉复形}。在有限域上，通过考察\term{frobenius}{Frobenius}在\term{cohomology}{上同调}各茎上的作用，可以定义\term{etale}{平展}\term{sheaf}{层}\term{complex}{复形}的\term{pure}{纯}性。Gabber的基本结果断言：\term{intersection-complex}{交叉复形}是\term{pure}{纯}的。将它与Deligne关于\term{pure}{纯}性在\term{direct-image}{正像}下行为的定理（Weil猜想）结合，便得到一个\term{decomposition-theorem}{分解定理}：\term{intersection-complex}{交叉复形}的\term{direct-image}{正像}可分解为带局部系数的\term{intersection-complex}{交叉复形}的直和。长篇论文[1]讨论的正是\term{perversity}{反常度}、\term{intersection-complex}{交叉复形}与\term{pure}{纯}性这一主题，也是本次会议的主要主题。

由此衍生的一个主题，是\term{intersection-homology}{交叉同调}在\term{nilpotent}{幂零}轨道研究[3]、奇异空间的Morse理论[6]及算术群商空间的$L^2$\term{cohomology}{上同调}[15]中的各种应用。

\SourcePage{9}
\term{intersection-homology}{交叉同调}似乎在所有情形下都带有Hodge结构，这应当与Gabber的\term{pure}{纯}性定理相呼应。一般情形尚无法证明，但[4]提出了一个应当给出Hodge\term{filtration}{滤过}的\term{filtration}{滤过}；[5]与[15]则说明，在某些情形下，利用$L^2$\term{cohomology}{上同调}的解释可以得到Hodge结构。

\term{intersection-homology}{交叉同调}上的这些Hodge结构，应当可以直接在与之对应的\term{holonomic-module}{完整模}上描述。[4]尝试完成这一点；对于对应于\term{vanishing-cycles}{消失轮换}的\term{holonomic-module}{完整模}，[11]及[13]已经给出了这样的描述。

与\term{vanishing-cycles}{消失轮换}对应的模还附带一个多项式，称为Bernstein多项式；对于孤立奇点，也称为$b$-函数。[9]将这个多项式与\term{vanishing-cycles}{消失轮换}空间上的\term{monodromy}{单值作用}联系起来。最后，其余报告围绕局部或整体Euler--Poincaré特征[10]、[7]、[2]，以及\term{vanishing-cycles}{消失轮换}[12]、[14]展开。
\begin{flushright}会议组织者\end{flushright}
\begin{center}
\begin{tabular}{@{}p{.46\textwidth}p{.46\textwidth}@{}}
B. Teissier & J. L. Verdier\\[5pt]
法国综合理工学院数学系 & 法国巴黎高等师范学校数学系\\
91128 Palaiseau Cedex & 45 rue d'Ulm, 75005 Paris
\end{tabular}
\end{center}
\UnnumberedChapter{论文集总目}
\SourcePage{10}
\section*{第1卷：Astérisque第100卷（1982）}
\noindent [1] A. A. Beilinson、J. Bernstein、P. Deligne：《\term{perverse-sheaf}{反常层}》（Faisceaux pervers）。
\begin{description}[font=\normalfont\bfseries]
\item[{引言}]
\item[{\term{triangulated-category}{三角范畴}的\term{abelian-subcategory}{阿贝尔子范畴}}] 1.1 \term{triangulated-category}{三角范畴}；1.2 \term{abelian-subcategory}{阿贝尔子范畴}；1.3 $t$-范畴；1.4 \term{recollement}{粘合}。
\item[{\term{stratification}{分层}空间与\term{scheme}{概形}上的\term{perverse-sheaf}{反常层}}] 2.1 \term{stratification}{分层}空间；2.2 \term{scheme}{概形}。
\item[{补充材料}] 3.1 \term{filtered-derived-category}{滤过导出范畴}、\term{canonical-filtration}{典范滤过}与\term{stupid-filtration}{朴素滤过}；3.2 \term{localization}{局部化}；3.3 整系数\term{cohomology}{上同调}。
\item[{\term{self-dual}{自对偶}\term{perversity}{反常度}：几何性质}] 4.1 \term{affine}{仿射}态射；4.2 正合性与\term{adjunction}{伴随}；4.3 \term{simple-object}{单对象}；4.4 \term{vanishing-cycles}{消失轮换}（上界估计）；4.5 \term{betti-number}{Betti数}的估计。
\item[{\term{self-dual}{自对偶}\term{perversity}{反常度}：\term{weight}{权}}] 5.1 回顾[1]；5.2 一个逆命题；5.3 \term{weight-filtration}{权滤过}；5.4 \term{pure}{纯}\term{complex}{复形}。
\item[{从有限域到复数域}] 6.1 原理；6.2 例子。
\item[{术语索引；参考文献}]
\end{description}
\SourcePage{11}
\section*{第2、3卷：Astérisque第101、102卷（1983）}
\begin{description}[leftmargin=2.8em,labelwidth=2.3em,font=\normalfont]
\item[{[2]}] J. P. Brasselet：双变量理论中\term{chern-class}{Chern类}的存在性（Existence des classes de Chern en théorie bivariante）。
\item[{[3]}] W. Borho、R. MacPherson：\term{nilpotent}{幂零}簇的部分消解（Partial resolutions of nilpotent varieties）。
\item[{[4]}] J. P. Brylinski：具有\term{regular-singularity}{正则奇点}的\term{holonomic-module}{完整模}与Hodge\term{filtration}{滤过}，II（Modules holonomes à singularités régulières et filtration de Hodge II）。
\item[{[5]}] J. Cheeger：复\term{cone}{锥}的Hodge理论（Hodge Theory of complex cones）。
\item[{[6]}] M. Goresky、R. MacPherson：Morse理论与\term{intersection-homology}{交叉同调}理论（Morse theory and Intersection Homology theory）。
\item[{[7]}] G. Laumon：曲面上\term{constructible}{可构造}\term{sheaf}{层}的Euler--Poincaré特征（Caractéristique d'Euler-Poincaré des faisceaux constructibles sur une surface）。
\item[{[8]}] G. Lusztig：奇点、特征标公式与\term{weight}{权}重数（Singularities, character formulas, weight multiplicities）。
\item[{[9]}] B. Malgrange：关于Boutet de Monvel与Kashiwara指标定理的报告（Rapport sur les théorèmes d'indice de Boutet de Monvel et Kashiwara）。
\item[{[10]}] B. Malgrange：Bernstein--Sato多项式与消失\term{cohomology}{上同调}（Polynômes de Bernstein-Sato et cohomologie évanescente）。
\item[{[11]}] F. Pham：具有孤立临界点的函数芽所附带的\term{mixed}{混合}Hodge结构（Structures de Hodge mixtes associées à un germe de fonction à point critique isolé）。
\item[{[12]}] C. Sabbah：无爆破的\term{stratification}{分层}解析态射与\term{vanishing-cycles}{消失轮换}（Morphismes analytiques stratifiés sans éclatement et cycles évanescents）。
\item[{[13]}] M. Saito：《Gauss--Manin系统与\term{mixed}{混合}Hodge结构》补遗（Supplement to ``Gauss-Manin system and mixed Hodge structure''）。
\item[{[14]}] J. L. Verdier：\term{sheaf}{层}的特殊化与温和\term{monodromy}{单值作用}（Spécialisation de faisceaux et monodromie modérée）。
\item[{[15]}] S. Zucker：Hodge理论与算术群（Hodge theory and arithmetic groups）。
\end{description}
\SourcePage{12}
\section*{1982年原版目录}
\noindent 下列页码为1982年原版印刷页码，予以保留供查阅；译本页码见本译本目录。
\begin{longtable}{@{}p{.85\textwidth}r@{}}
\toprule
0. 引言&7\\
1. \term{triangulated-category}{三角范畴}的\term{abelian-subcategory}{阿贝尔子范畴}&18\\
\quad 1.1 \term{triangulated-category}{三角范畴}；1.2 \term{abelian-subcategory}{阿贝尔子范畴}；1.3 $t$-范畴；1.4 \term{recollement}{粘合}&\\
2. \term{stratification}{分层}空间与\term{scheme}{概形}上的\term{perverse-sheaf}{反常层}&56\\
\quad 2.1 \term{stratification}{分层}空间；2.2 \term{scheme}{概形}&\\
3. 补充材料&76\\
\quad 3.1 \term{filtered-derived-category}{滤过导出范畴}、\term{canonical-filtration}{典范滤过}与\term{stupid-filtration}{朴素滤过}；3.2 \term{localization}{局部化}；3.3 整系数\term{cohomology}{上同调}&\\
4. \term{self-dual}{自对偶}\term{perversity}{反常度}：几何性质&101\\
\quad 4.1 \term{affine}{仿射}态射；4.2 正合性与\term{adjunction}{伴随}；4.3 \term{simple-object}{单对象}；4.4 \term{vanishing-cycles}{消失轮换}（上界估计）；4.5 \term{betti-number}{Betti数}的估计&\\
5. \term{self-dual}{自对偶}\term{perversity}{反常度}：\term{weight}{权}&122\\
\quad 5.1 回顾[1]；5.2 一个逆命题；5.3 \term{weight-filtration}{权滤过}；5.4 \term{pure}{纯}\term{complex}{复形}&\\
6. 从有限域到复数域&148\\
\quad 6.1 原理；6.2 例子&\\
参考文献&167\\
术语索引&169\\
符号索引&171\\\bottomrule
\end{longtable}
\SourcePage{13}
\section*{主题分类}
\begin{enumerate}[label=\alph*)]
\item \term{perversity}{反常度}、\term{intersection-complex}{交叉复形}、\term{pure}{纯}性：[1]。
\item 交叉理论及其应用：[1]、[3]、[6]、[8]、[15]。
\item 交叉理论与Hodge结构：[4]、[5]、[15]。
\item $\mathrm R\Psi$上的Hodge结构：[11]、[13]。
\item $\mathrm R\Psi$与$\chi$：[9]、[14]、[2]、[7]、[12]、[10]。
\end{enumerate}
